% ============================================================================
% 05_reinforcement_learning/01_theory.tex
% Section body only: the master manuscript supplies the chapter/section title.
% Figures are source files in the project-root figures/ directory and are
% included with \input so that the diagrams remain editable LaTeX material.
% ============================================================================

The preceding discussion of model architecture, optimization, and scaling
establishes how a language model can be represented and trained as a predictive
system. It does not, by itself, specify how the model should change its behavior
when the quality of a generated sequence is evaluated through a reward signal.
That problem introduces a different mathematical object: a decision-making
policy whose actions affect not only the immediate outcome but also the states
and decisions that follow. Reinforcement learning provides a framework for
reasoning about this sequential dependence.

The purpose of this section is to establish the mathematical foundations needed
to understand reinforcement learning as an optimization problem over policies.
The development begins with Markov decision processes and Bellman operators,
then proceeds to policy improvement and approximation error, information
geometry and natural gradients, trust-region optimization, and regularized
Markov decision processes. These topics are connected by a common question:
how can a policy be improved while controlling the consequences of imperfect
value estimates, changes in the distribution of visited states, and finite
optimization steps? The cited papers provide the main theoretical structure
for this development; the final connection to autoregressive language models
is an application of that structure to the subject of this book.

\subsection{Markov decision processes, policies, and value functions}
\label{subsec:rl-mdp-value-functions}

A reinforcement-learning problem is a sequential decision problem. At each
decision time, an agent receives information about the current situation,
selects an action, and observes the consequences of that action. The next
situation depends on the current situation and action, and the reward may be
delayed. Thus, maximizing the reward at the current time is generally not
equivalent to maximizing the total reward accumulated over an interaction.
The mathematical formulation used throughout the cited work is the discounted
Markov decision process (MDP).

For the finite-state and finite-action setting used to develop the matrix
identities, an MDP is specified by
\begin{equation}
    \mathcal{M}
    = (\mathcal{S},\mathcal{A},P,r,\gamma,\rho_0),
    \label{eq:rl-mdp-tuple}
\end{equation}
where $\mathcal{S}$ is the state space, $\mathcal{A}$ is the action space,
$P(s'\mid s,a)$ is the transition probability, $r(s,a)$ is the expected
immediate reward, $\gamma\in(0,1)$ is the discount factor, and $\rho_0$ is the
distribution of the initial state. For every state--action pair,
\begin{equation}
    P(s'\mid s,a)\geq 0,
    \qquad
    \sum_{s'\in\mathcal{S}}P(s'\mid s,a)=1.
    \label{eq:rl-transition-normalization}
\end{equation}
The finite-state assumption is used only when vectors and matrices are written
explicitly. The same definitions extend to general state spaces by replacing
finite sums with integrals and transition matrices with probability kernels,
as in the formulation considered by \textcite{scherrer2014approximate}.

At time $t$, the agent is in state $s_t$, chooses action $a_t$, receives reward
$r_t$, and the environment produces state $s_{t+1}$. The transition model is

\begin{equation}
    P(s'\mid s,a)
    = \Pr(s_{t+1}=s'\mid s_t=s,a_t=a).
    \label{eq:rl-transition-kernel}
\end{equation}

The Markov property states that, conditioned on the current state and action,
the next-state distribution does not depend on the preceding history:
\begin{equation}
\begin{split}
    &\Pr(s_{t+1}\mid s_t,a_t,s_{t-1},a_{t-1},\ldots,s_0)\\
    &\hspace{3em}=P(s_{t+1}\mid s_t,a_t).
\end{split}
\label{eq:rl-markov-property}
\end{equation}
This is a conditional-independence assumption, not an assertion that the
environment is deterministic. The transition may remain stochastic. The
assumption says that the state representation contains the information needed
to determine the conditional distribution of the next state. If the chosen
state representation omits information relevant to future transitions or
rewards, the process represented by that state need not be Markov.

The policy specifies how the agent selects actions. A stochastic stationary
policy is a conditional probability distribution
\begin{equation}
    \pi(a\mid s)
    =\Pr(a_t=a\mid s_t=s),
    \label{eq:rl-policy}
\end{equation}
with $\pi(\cdot\mid s)$ belonging to the probability simplex
\begin{equation}
    \Delta_{\mathcal{A}}
    =\left\{p\in\mathbb{R}^{|\mathcal{A}|}:
    p(a)\geq0,\ \sum_{a\in\mathcal{A}}p(a)=1\right\}.
    \label{eq:rl-probability-simplex}
\end{equation}
A deterministic policy is a special case in which the distribution assigns
probability one to a single action. The stochastic formulation is more general
and is necessary for the probability-distribution optimization developed
later in this section.

\begin{figure}[htbp]
    \centering
    \resizebox{\textwidth}{!}{% Native LaTeX picture: no external graphics package is required.
\setlength{\unitlength}{1mm}
\begin{picture}(150,48)
    \put(8,22){\framebox(48,18){\shortstack{Agent\\policy $\pi_\theta(a\mid s)$}}}
    \put(94,22){\framebox(48,18){\shortstack{Environment\\$P(s'\mid s,a),\ r(s,a)$}}}
    \put(56,34){\vector(1,0){38}}
    \put(61,38){\makebox(30,4){action $a_t$}}
    \put(94,27){\vector(-1,0){38}}
    \put(61,17){\makebox(32,4){$s_{t+1},\ r_t$}}
    \put(8,9){\makebox(48,5){observation $s_t$}}
    \put(94,9){\makebox(48,5){transition and reward}}
\end{picture}
}
    \caption{The interaction represented by a Markov decision process. The agent
    selects an action from its policy, and the environment returns the next
    state and reward according to its transition and reward model.}
    \label{fig:rl-mdp-loop}
\end{figure}
% Alt text for accessibility register: Two boxes represent the agent and the
% environment. An action arrow points from the agent to the environment; a
% state-and-reward arrow points back to the agent.

The objective must account for the complete future reward sequence. For an
infinite-horizon discounted problem, the return beginning at time $t$ is
\begin{equation}
    G_t=\sum_{k=0}^{\infty}\gamma^k r_{t+k}.
    \label{eq:rl-return}
\end{equation}
The discount factor assigns weight one to the immediate reward and weight
$\gamma^k$ to a reward received $k$ steps later. If rewards are bounded,
$|r_t|\leq R_{\max}$, then
\begin{equation}
    |G_t|
    \leq R_{\max}\sum_{k=0}^{\infty}\gamma^k
    =\frac{R_{\max}}{1-\gamma}.
    \label{eq:rl-return-bound}
\end{equation}
The geometric-series identity requires $0<\gamma<1$ and guarantees a bounded
return under the bounded-reward assumption. As $\gamma$ approaches one, more
distant rewards receive greater relative weight and the bound
$R_{\max}/(1-\gamma)$ increases.

The state-value function of a policy is the conditional expected return when
the process starts at state $s$ and subsequently follows that policy:
\begin{equation}
    V^\pi(s)
    =\mathbb{E}_{\pi}\!\left[
    \sum_{k=0}^{\infty}\gamma^k r_{t+k}
    \;\middle|\;s_t=s\right].
    \label{eq:rl-state-value}
\end{equation}
The expectation is over both the policy's action distribution and the
environment's transition distribution. The state--action value function
conditions additionally on the first action:
\begin{equation}
    Q^\pi(s,a)
    =\mathbb{E}\!\left[
    \sum_{k=0}^{\infty}\gamma^k r_{t+k}
    \;\middle|\;s_t=s,\ a_t=a,\ a_{t+j}\sim\pi\ (j\geq1)
    \right].
    \label{eq:rl-action-value}
\end{equation}
Thus, $V^\pi(s)$ evaluates the policy's action distribution at the current
state, whereas $Q^\pi(s,a)$ evaluates a specified first action followed by the
policy thereafter. Their relationship is
\begin{equation}
    V^\pi(s)
    =\sum_{a\in\mathcal{A}}\pi(a\mid s)Q^\pi(s,a).
    \label{eq:rl-vq-relation}
\end{equation}

The advantage function compares an action with the average value of the current
policy at the same state:
\begin{equation}
    A^\pi(s,a)=Q^\pi(s,a)-V^\pi(s).
    \label{eq:rl-advantage}
\end{equation}
A positive advantage means that the specified action has a greater expected
return than the policy average at that state; a negative advantage means that
its expected return is lower. From \cref{eq:rl-vq-relation} and the
normalization of the policy,
\begin{align}
    \sum_a\pi(a\mid s)A^\pi(s,a)
    &=\sum_a\pi(a\mid s)Q^\pi(s,a)
      -V^\pi(s)\sum_a\pi(a\mid s)\\
    &=V^\pi(s)-V^\pi(s)=0.
    \label{eq:rl-zero-mean-advantage}
\end{align}
The zero-mean identity is central to policy optimization: the expected
advantage under the current policy is zero at every state. The performance
analysis of trust-region methods uses this identity when constructing a local
surrogate objective \parencite{schulman2015trpo}.

\subsection{Bellman operators and dynamic programming}
\label{subsec:rl-bellman-dynamic-programming}

The value functions are not only expectations over trajectories; they satisfy
recursive equations. Decomposing the return into the immediate reward and the
discounted return from the next time step gives
\begin{equation}
    G_t=r_t+\gamma G_{t+1}.
    \label{eq:rl-return-recursion}
\end{equation}
Taking conditional expectations and using the Markov property yields the
Bellman expectation equation:
\begin{equation}
    V^\pi(s)
    =\sum_a\pi(a\mid s)
    \left[r(s,a)+\gamma\sum_{s'}P(s'\mid s,a)V^\pi(s')\right].
    \label{eq:rl-bellman-expectation}
\end{equation}
It states that the value of a state is the expected immediate reward plus the
discounted value of the next state.

For a finite state space with $n=|\mathcal{S}|$ and $m=|\mathcal{A}|$, define
the policy reward vector $r_\pi\in\mathbb{R}^{n}$ and transition matrix
$P_\pi\in\mathbb{R}^{n\times n}$ by
\begin{equation}
    r_\pi(s)=\sum_a\pi(a\mid s)r(s,a),
    \qquad
    P_\pi(s'\mid s)=\sum_a\pi(a\mid s)P(s'\mid s,a).
    \label{eq:rl-policy-induced-model}
\end{equation}
Each row of $P_\pi$ is a probability distribution. Hence $P_\pi$ is
row-stochastic, has nonnegative entries, and satisfies
$\|P_\pi\|_\infty=1$ and spectral radius $\rho(P_\pi)\leq1$. In vector form,
\cref{eq:rl-bellman-expectation} becomes
\begin{equation}
    V^\pi=r_\pi+\gamma P_\pi V^\pi.
    \label{eq:rl-bellman-vector}
\end{equation}
Rearranging gives
\begin{equation}
    (I-\gamma P_\pi)V^\pi=r_\pi,
    \qquad
    V^\pi=(I-\gamma P_\pi)^{-1}r_\pi.
    \label{eq:rl-policy-value-linear-system}
\end{equation}
The inverse exists because
$\rho(\gamma P_\pi)\leq\gamma<1$. Equivalently, the Neumann series converges:
\begin{equation}
    (I-\gamma P_\pi)^{-1}
    =\sum_{k=0}^{\infty}(\gamma P_\pi)^k,
    \qquad
    V^\pi=\sum_{k=0}^{\infty}\gamma^kP_\pi^k r_\pi.
    \label{eq:rl-neumann-series}
\end{equation}
The term $P_\pi^k r_\pi$ is the expected reward vector after $k$ transitions
under the policy. The Neumann expansion therefore recovers the discounted
trajectory interpretation of the value function through linear algebra.
The inverse in \cref{eq:rl-policy-value-linear-system} is a theoretical characterization, not a recommendation to form a dense matrix inverse numerically. A practical implementation solves the linear system with a sparse solver or uses iterative Bellman evaluation. The bound
\begin{equation}
    \|(I-\gamma P_\pi)^{-1}\|_\infty
    \leq\sum_{k=0}^{\infty}\gamma^k
    =\frac{1}{1-\gamma}
    \label{eq:rl-policy-evaluation-conditioning}
\end{equation}
shows why the discounted system can become more sensitive as $\gamma$ approaches one: the inverse norm bound grows with the effective horizon. Iterative evaluation remains contractive, but its worst-case contraction factor is $\gamma$, so the number of iterations required for a fixed error tolerance also increases as $\gamma$ approaches one.

For an arbitrary candidate value vector $v\in\mathbb{R}^{n}$, define the
Bellman evaluation operator
\begin{equation}
    [T_\pi v](s)
    =\sum_a\pi(a\mid s)
    \left[r(s,a)+\gamma\sum_{s'}P(s'\mid s,a)v(s')\right].
    \label{eq:rl-bellman-evaluation-operator}
\end{equation}
The value function is its fixed point, $V^\pi=T_\pi V^\pi$. The operator is a
contraction in the supremum norm. For any $v_1,v_2$,
\begin{align}
    \|T_\pi v_1-T_\pi v_2\|_\infty
    &=\gamma\|P_\pi(v_1-v_2)\|_\infty\\
    &\leq\gamma\|P_\pi\|_\infty\|v_1-v_2\|_\infty\\
    &=\gamma\|v_1-v_2\|_\infty.
    \label{eq:rl-policy-contraction}
\end{align}
Because $\gamma<1$, repeated evaluation $v_{k+1}=T_\pi v_k$ converges to the
unique fixed point $V^\pi$ for any initial vector $v_0$.

Policy evaluation alone does not optimize the policy. To express optimization,
replace the action average under a fixed policy with a maximum over actions.
The Bellman optimality operator is
\begin{equation}
    [T^*v](s)
    =\max_{a\in\mathcal{A}}
    \left[r(s,a)+\gamma\sum_{s'}P(s'\mid s,a)v(s')\right].
    \label{eq:rl-bellman-optimality}
\end{equation}
Its unique fixed point is the optimal value function $V^*$, and an optimal
deterministic policy is greedy with respect to $V^*$:
\begin{equation}
    V^*=T^*V^*,
    \qquad
    \pi^*(s)\in\operatorname*{arg\,max}_{a}
    \left[r(s,a)+\gamma\sum_{s'}P(s'\mid s,a)V^*(s')\right].
    \label{eq:rl-optimal-policy}
\end{equation}
The inequality
$|\max_a x_a-\max_a y_a|\leq\max_a|x_a-y_a|$ together with the stochastic
normalization of $P$ gives
\begin{equation}
    \|T^*v_1-T^*v_2\|_\infty
    \leq\gamma\|v_1-v_2\|_\infty.
    \label{eq:rl-optimality-contraction}
\end{equation}
Thus, value iteration $v_{k+1}=T^*v_k$ converges to $V^*$.

Policy iteration organizes the same fixed-point problem differently. It first
evaluates a current policy, then improves it by choosing actions that maximize
the one-step return computed from the current value function. If
$G(v)=\operatorname*{arg\,max}_{\pi}T_\pi v$ denotes the greedy policy set,
the exact update is
\begin{equation}
    v_k=V^{\pi_k},
    \qquad
    \pi_{k+1}\in G(v_k).
    \label{eq:rl-policy-iteration}
\end{equation}
Modified policy iteration interpolates between a single optimality update and
complete policy evaluation:
\begin{equation}
    \pi_{k+1}\in G(v_k),
    \qquad
    v_{k+1}=(T_{\pi_{k+1}})^m v_k.
    \label{eq:rl-modified-policy-iteration}
\end{equation}
For $m=1$, the evaluation and greedy steps combine into a value-iteration
update. As $m$ increases, the new policy is evaluated more fully; in the limit
of exact evaluation, the update becomes policy iteration. This continuum is the
starting point for the approximate and regularized schemes considered below
\parencite{scherrer2014approximate,geist2019regularized}.

\begin{figure}[htbp]
    \centering
    \resizebox{\textwidth}{!}{% Native LaTeX picture: value estimate -> approximate greedy policy -> evaluation.
\setlength{\unitlength}{1mm}
\begin{picture}(155,55)
    \put(3,35){\framebox(34,13){\shortstack{Value estimate\\$v_k$}}}
    \put(49,35){\framebox(46,13){\shortstack{Approximate greedy\\$G_{\varepsilon}$}}}
    \put(108,35){\framebox(40,13){\shortstack{New policy\\$\pi_{k+1}$}}}
    \put(94,5){\framebox(54,13){\shortstack{Policy evaluation\\$(T_{\pi_{k+1}})^m$}}}
    \put(37,41){\vector(1,0){12}}
    \put(95,41){\vector(1,0){13}}
    \put(128,35){\vector(0,-1){17}}
    \put(94,11){\vector(-1,0){53}}
    \put(51,14){\makebox(37,4){updated value $v_{k+1}$}}
\end{picture}
}
    \caption{The modified policy-iteration cycle. A value estimate is passed to
    an approximate greedy step, the resulting policy is evaluated for $m$
    Bellman steps, and the updated value estimate begins the next iteration.}
    \label{fig:rl-policy-iteration}
\end{figure}
% Alt text for accessibility register: A left-to-right approximate greedy
% update is followed by a downward policy-evaluation step and a return to the
% updated value estimate.

\subsection{Policy improvement, performance difference, and approximation}
\label{subsec:rl-policy-improvement-approximation}

The Bellman optimality equation characterizes an optimal value function, but
policy optimization also requires a relation between the performance of two
policies. Let $\rho_0$ be the initial-state distribution and define
\begin{equation}
    \eta(\pi)=\mathbb{E}_{s_0\sim\rho_0}[V^\pi(s_0)].
    \label{eq:rl-policy-performance}
\end{equation}
For a policy $\pi$, its unnormalized discounted state-visitation frequency is
\begin{equation}
    \rho_\pi(s)
    =\sum_{t=0}^{\infty}\gamma^t
      \Pr(s_t=s\mid s_0\sim\rho_0,\pi).
    \label{eq:rl-discounted-occupancy}
\end{equation}
The normalized discounted occupancy distribution is
$d_{\pi,\rho_0}=(1-\gamma)\rho_\pi$ and, in row-vector notation,
\begin{equation}
    d_{\pi,\rho_0}
    =(1-\gamma)\rho_0(I-\gamma P_\pi)^{-1}.
    \label{eq:rl-normalized-occupancy}
\end{equation}
The normalization follows from
$(1-\gamma)\sum_{t=0}^{\infty}\gamma^t=1$. This distribution weights states
according to how often they are encountered under the policy, with later
visits discounted.

For two policies $\pi$ and $\widetilde\pi$, the performance difference
identity is
\begin{equation}
    \eta(\widetilde\pi)
    =\eta(\pi)
    +\sum_s\rho_{\widetilde\pi}(s)
      \sum_a\widetilde\pi(a\mid s)A^\pi(s,a).
    \label{eq:rl-performance-difference}
\end{equation}
To see the source of the identity, write
$A^\pi(s_t,a_t)=r_t+\gamma V^\pi(s_{t+1})-V^\pi(s_t)$ in conditional expectation,
then sum the discounted terms along a trajectory generated by $\widetilde\pi$.
The value terms telescope, leaving the difference between the new and old
expected returns. Equation~\cref{eq:rl-performance-difference} shows that
policy improvement depends on two coupled quantities: the advantages of the
new policy's actions and the state-visitation distribution induced by the new
policy. If the expected advantage is nonnegative at every state, the new policy
cannot have lower performance. Exact greedy improvement satisfies this
condition; approximate improvement need not.

In approximate policy iteration, the exact greedy operator is replaced by an
approximate one. For a reference distribution $\nu$, an approximate greedy
policy can be characterized by
\begin{equation}
    \nu(Tv-T_\pi v)\leq\varepsilon,
    \qquad
    \nu f:=\mathbb{E}_{s\sim\nu}[f(s)],
    \label{eq:rl-approx-greedy}
\end{equation}
where $\varepsilon\geq0$ is the greedy-step error. This is an average error
under $\nu$, not a uniform guarantee over every state. In the exact case,
$\varepsilon=0$, the policy is greedy on the support of $\nu$. The distinction
matters because the distribution used to fit or evaluate the approximate greedy
step need not match the state distribution induced by the resulting policy.

The mismatch can be quantified using concentrability coefficients. Given a
evaluation distribution $\mu$ and a reference distribution $\nu$, let $c(i)$ be
a coefficient satisfying, componentwise,
\begin{equation}
    \mu P_{\pi_1}P_{\pi_2}\cdots P_{\pi_i}
    \leq c(i)\nu
    \label{eq:rl-concentrability}
\end{equation}
for all admissible sequences of policies of length $i$. A small $c(i)$ means
that distributions reached after $i$ transitions remain covered by the
reference distribution. A large or infinite coefficient indicates that the
evaluation distribution can place substantial mass on states that are poorly
represented under $\nu$. Weighted sums of the coefficients, with powers of
$\gamma$, appear in performance bounds because errors propagated through many
transitions are discounted but are not eliminated. Thus, an approximation error
alone does not characterize the quality of an approximate policy-iteration
method; the distribution mismatch appearing in its bound is also essential
\parencite{scherrer2014approximate}.

The standard Approximate Policy Iteration update is
\begin{equation}
    \pi_{k+1}
    =G_{\varepsilon_{k+1}}(\nu,V^{\pi_k}),
    \label{eq:rl-api}
\end{equation}
where the notation makes explicit that the approximate greedy operator depends
on both the value estimate and the distribution under which its error is
measured. Conservative Policy Iteration controls the change of policy by
mixing the current policy with an approximately greedy policy:
\begin{equation}
    \pi_{k+1}
    =(1-\alpha_{k+1})\pi_k
      +\alpha_{k+1}
      G_{\varepsilon_{k+1}}(d_{\pi_k,\rho_0},V^{\pi_k}).
    \label{eq:rl-cpi}
\end{equation}
Here $\alpha_{k+1}\in(0,1]$ is a step size. The occupancy distribution is used
in the greedy step, and a sufficiently conservative step can be selected to
obtain an improvement guarantee. The API$(\alpha)$ variant also mixes policies but uses the reference
distribution $\nu$ in its approximate greedy step:
\begin{equation}
    \pi_{k+1}=(1-\alpha)\pi_k
       +\alpha G_{\varepsilon_{k+1}}(\nu,V^{\pi_k}).
    \label{eq:rl-api-alpha}
\end{equation}
This is simpler to implement when sampling from the current policy's
occupancy distribution is inconvenient.

A different construction retains a sequence of policies rather than replacing
the current policy by a mixture. Given stationary policies
$\pi_1,\ldots,\pi_k$, define the finite non-stationary sequence
$\sigma_k=\pi_k\pi_{k-1}\cdots\pi_1$, where the first action uses $\pi_k$, the
next uses $\pi_{k-1}$, and so forth. Its value is
\begin{equation}
    V^{\sigma_k}
    =T_{\pi_k}T_{\pi_{k-1}}\cdots T_{\pi_1}r.
    \label{eq:rl-psdp-value}
\end{equation}
The infinite-horizon extension called $\mathrm{PSDP}_\infty$ constructs the
sequence iteratively through
\begin{equation}
    \pi_{k+1}
    =G_{\varepsilon_{k+1}}(\nu,V^{\sigma_k}),
    \label{eq:rl-psdp}
\end{equation}
and can return a periodic policy that repeats the retained finite sequence.
The Non-Stationary Policy Iteration method with period $m$ retains only the
most recent $m$ policies and repeatedly executes
$\sigma_k^m=\pi_k\pi_{k-1}\cdots\pi_{k-m+1}$. Its next policy is obtained
from the value of the periodic extension of this sequence:
\begin{equation}
    \pi_{k+1}=G_{\varepsilon_{k+1}}
       \bigl(\nu,V^{(\sigma_k^m)^\infty}\bigr).
    \label{eq:rl-nspi}
\end{equation}
It therefore creates a continuum between the one-policy memory of API and
the longer memory required by $\mathrm{PSDP}_\infty$.

\begin{table}[htbp]
    \centering
    \caption{Qualitative comparison of the approximate policy-iteration schemes
    analyzed by \textcite{scherrer2014approximate}. The guarantees depend on the
    concentrability coefficients and approximation assumptions, not only on
    the algorithm name.}
    \label{tab:rl-approximate-policy-iteration}
    \begin{tabularx}{\textwidth}{lXX}
        \toprule
        Scheme & Policy update structure & Principal theoretical or computational trade-off \\
        \midrule
        API & Replace the current policy with an approximately greedy policy. & Constant memory; its bounds use more restrictive concentrability coefficients. \\
        CPI & Mix the current policy with an approximately greedy policy computed under its occupancy distribution. & Can have stronger performance guarantees, but may require substantially more iterations and occupancy-distribution sampling. \\
        $\mathrm{PSDP}_\infty$ & Improve a growing non-stationary sequence of policies. & Combines a strong guarantee with an API-like iteration requirement, at the cost of memory growing with the number of iterations. \\
        NSPI$(m)$ & Repeat a sequence containing only the most recent $m$ policies. & Makes memory proportional to $m$ and exposes a trade-off between memory and the performance bound. \\
        \bottomrule
    \end{tabularx}
\end{table}

The comparison in \cref{tab:rl-approximate-policy-iteration} should not be read
as a universal ranking independent of the MDP. Scherrer's analysis shows that
the hierarchy of concentrability coefficients can make one guarantee finite
while another is infinite. Under the bounds analyzed in that paper, CPI can
have a substantially stronger guarantee than API, but its iteration requirement
can grow on the order of $1/\varepsilon^2$, whereas API has a logarithmic
iteration dependence on $1/\varepsilon$. The $\mathrm{PSDP}_\infty$ construction
combines a CPI-like guarantee with an API-like iteration requirement, while
requiring storage for the policies accumulated along the sequence. NSPI$(m)$
limits that storage to a fixed window. The conclusion is methodological:
performance, iteration complexity, memory, and state-distribution coverage must
be analyzed together \parencite{scherrer2014approximate}.

\subsection{Information geometry and covariant policy search}
\label{subsec:rl-information-geometry}

When a policy is represented by a differentiable model, its probability
function depends on a parameter vector $\theta\in\mathbb{R}^d$:
$\pi_\theta(a\mid s)$. The objective is to maximize the expected return
$J(\theta)=\eta(\pi_\theta)$. Ordinary gradient ascent takes the form
\begin{equation}
    \theta_{k+1}
    =\theta_k+\alpha\nabla_\theta J(\theta_k).
    \label{eq:rl-euclidean-gradient}
\end{equation}
This update measures step size using Euclidean distance in parameter space.
However, the parameter vector is a coordinate description of a policy, not the
behavioral object itself. Two different parameterizations can represent the
same family of probability distributions while assigning different Euclidean
lengths and directions to corresponding changes in policy behavior.

Suppose that another parameterization is defined by a differentiable invertible
coordinate transformation $\phi=f(\theta)$. For a small displacement,
$\Delta\phi\simeq J_f(\theta)\Delta\theta$, where $J_f$ is the Jacobian of the
coordinate transformation. A geometric update should transform consistently
with this change of coordinates. An ordinary Euclidean gradient generally does
not have this property. This is the non-covariance problem addressed by
\textcite{bagnell2003covariant}.

A Riemannian metric assigns an inner product to local tangent directions. In
parameter coordinates, let $F(\theta)$ be a positive-definite metric tensor on
the identifiable tangent directions. The local squared length of an update is
\begin{equation}
    \|\Delta\theta\|_{F}^{2}
    =\Delta\theta^{\mathsf T}F(\theta)\Delta\theta.
    \label{eq:rl-riemannian-length}
\end{equation}
The direction of steepest first-order ascent under a fixed metric length is the
solution of
\begin{equation}
    \max_{\Delta\theta}\quad
    \nabla_\theta J(\theta)^{\mathsf T}\Delta\theta
    \quad\text{subject to}\quad
    \Delta\theta^{\mathsf T}F(\theta)\Delta\theta\leq\delta^2.
    \label{eq:rl-natural-gradient-problem}
\end{equation}
Introducing a Lagrange multiplier and differentiating with respect to
$\Delta\theta$ gives
$\nabla_\theta J-\lambda F\Delta\theta=0$. Therefore, up to a positive step
length, the natural-gradient direction is
\begin{equation}
    \widetilde{\nabla}_\theta J(\theta)
    =F(\theta)^{-1}\nabla_\theta J(\theta).
    \label{eq:rl-natural-gradient}
\end{equation}
If the metric is singular because the parameterization contains redundant
directions, the inverse is understood on the identifiable subspace or replaced
by a suitable generalized inverse. The metric changes the direction, not merely
the scalar learning rate.

The metric used by covariant policy search is derived from the probability
distribution over trajectories induced by the stochastic controller. For a
finite trajectory $\tau=(s_0,a_0,\ldots,s_T)$, its probability is
\begin{equation}
    p_\theta(\tau)
    =\rho_0(s_0)
      \prod_{t=0}^{T-1}
      \pi_\theta(a_t\mid s_t)P(s_{t+1}\mid s_t,a_t).
    \label{eq:rl-trajectory-probability}
\end{equation}
When the initial-state distribution and environment dynamics do not depend on
$\theta$, the trajectory score is
\begin{equation}
    \nabla_\theta\log p_\theta(\tau)
    =\sum_{t=0}^{T-1}
      \nabla_\theta\log\pi_\theta(a_t\mid s_t).
    \label{eq:rl-trajectory-score}
\end{equation}
The Fisher information matrix of the trajectory distribution is
\begin{equation}
    F_T(\theta)
    =\mathbb{E}_{\tau\sim p_\theta}
      \left[
      \nabla_\theta\log p_\theta(\tau)
      \nabla_\theta\log p_\theta(\tau)^{\mathsf T}
      \right].
    \label{eq:rl-trajectory-fisher}
\end{equation}
It is positive semidefinite because, for every vector $x$,
$x^{\mathsf T}F_Tx$ is the expectation of the nonnegative random variable
$(x^{\mathsf T}\nabla_\theta\log p_\theta(\tau))^2$. The Fisher matrix measures
how sensitively the trajectory distribution changes when the policy parameters
change.
The same trajectory score gives the likelihood-ratio policy-gradient identity. For an episodic trajectory return $R(\tau)$,
\begin{align}
    \nabla_\theta J(\theta)
    &=\nabla_\theta\int p_\theta(\tau)R(\tau)\,d\tau\\
    &=\int p_\theta(\tau)R(\tau)
      \nabla_\theta\log p_\theta(\tau)\,d\tau\\
    &=\mathbb{E}_{\tau\sim p_\theta}
      \left[R(\tau)\nabla_\theta\log p_\theta(\tau)\right].
    \label{eq:rl-score-function-gradient}
\end{align}
The second equality follows from $\nabla p=p\nabla\log p$ and requires that differentiation can be interchanged with integration. Substituting the trajectory score from \cref{eq:rl-trajectory-score} expresses the gradient as a sum of per-decision score terms. A state-dependent baseline can be subtracted without changing the expectation because $\mathbb{E}_{a\sim\pi_\theta(\cdot\mid s)}[\nabla_\theta\log\pi_\theta(a\mid s)]=0$. Choosing a value baseline gives the advantage-weighted form used in policy optimization:
\begin{equation}
    \nabla_\theta J(\theta)
    =\mathbb{E}_{\pi_\theta}
      \left[\sum_t \gamma^t\nabla_\theta\log\pi_\theta(a_t\mid s_t)
      A^{\pi_\theta}(s_t,a_t)\right],
    \label{eq:rl-advantage-policy-gradient}
\end{equation}
with the trajectory weighting and discount convention chosen consistently with $J$. This identity explains why advantage estimation is central to the actor--critic methods developed in the next section.

For the discounted start-state formulation, the corresponding state weighting
is the normalized discounted occupancy distribution. The per-decision policy
Fisher matrix can consequently be written as
\begin{equation}
    F(\theta)
    =\mathbb{E}_{\substack{s\sim d_{\pi_\theta,\rho_0}\\
                            a\sim\pi_\theta(\cdot\mid s)}}
      \left[
      \nabla_\theta\log\pi_\theta(a\mid s)
      \nabla_\theta\log\pi_\theta(a\mid s)^{\mathsf T}
      \right],
    \label{eq:rl-occupancy-fisher}
\end{equation}
up to the normalization convention used for the trajectory objective. In an
ergodic average-reward formulation, the analogous weighting is the stationary
state distribution. This construction makes the metric depend on the
probabilistic behavior of the controller rather than on an arbitrary Euclidean
coordinate system \parencite{bagnell2003covariant}.

The Fisher metric is locally related to the Kullback--Leibler divergence. For
nearby parameters and a regular differentiable model,
\begin{equation}
    D_{\mathrm{KL}}(p_\theta\|p_{\theta+\Delta\theta})
    =\frac{1}{2}\Delta\theta^{\mathsf T}F_T(\theta)
      \Delta\theta+o(\|\Delta\theta\|^2).
    \label{eq:rl-fisher-kl-expansion}
\end{equation}
Thus, the Fisher matrix is the local quadratic form of statistical
 distinguishability between nearby trajectory distributions. This explains why
the natural gradient is a steepest-ascent direction when policy changes are
measured through their effect on probability distributions.

The same geometry connects to compatible actor--critic approximation. Define
the policy score feature
$\psi_\theta(s,a)=\nabla_\theta\log\pi_\theta(a\mid s)$ and a compatible
advantage approximator
\begin{equation}
    A_w(s,a)=w^{\mathsf T}\psi_\theta(s,a).
    \label{eq:rl-compatible-critic}
\end{equation}
Under the same state--action weighting used to define the policy gradient,
the least-squares normal equations for this compatible approximation take the
form $F(\theta)w=g$, where $g=\nabla_\theta J(\theta)$, provided the compatible
features and the policy-gradient identity are used consistently. Hence
\begin{equation}
    w=F(\theta)^{-1}g.
    \label{eq:rl-compatible-natural-gradient}
\end{equation}
This result explains the relationship between the geometry of the policy
family and actor--critic estimation discussed by \textcite{bagnell2003covariant}.
It does not imply that an arbitrary critic architecture is compatible; the
feature form in \cref{eq:rl-compatible-critic} is the relevant condition.

\begin{figure}[htbp]
    \centering
    \resizebox{\textwidth}{!}{% Native LaTeX picture: parameterization, path distribution, and metric update.
\setlength{\unitlength}{1mm}
\begin{picture}(160,58)
    \put(2,38){\framebox(36,13){\shortstack{Parameters\\$\theta\in\mathbb{R}^d$}}}
    \put(62,38){\framebox(36,13){\shortstack{Policy\\$\pi_\theta(a\mid s)$}}}
    \put(122,38){\framebox(36,13){\shortstack{Path distribution\\$p_\theta(\tau)$}}}
    \put(38,44){\vector(1,0){24}}
    \put(98,44){\vector(1,0){24}}
    \put(122,15){\framebox(36,13){\shortstack{Fisher metric\\$F(\theta)$}}}
    \put(62,15){\framebox(36,13){\shortstack{Natural direction\\$F^{-1}g$}}}
    \put(2,15){\framebox(36,13){\shortstack{Trust region\\$D_{\rm KL}\leq\delta$}}}
    \put(140,38){\vector(0,-1){10}}
    \put(122,21){\vector(-1,0){24}}
    \put(62,21){\vector(-1,0){24}}
    \put(43,7){\makebox(75,5){local probability geometry controls the update}}
\end{picture}
}
    \caption{The information-geometric route from policy parameters to a
    trajectory distribution, its Fisher metric, and a trust-region-limited
    natural-gradient update.}
    \label{fig:rl-information-geometry}
\end{figure}
% Alt text for accessibility register: Parameter coordinates define a policy
% and a trajectory distribution. The trajectory distribution defines a Fisher
% metric, which is used to scale the policy gradient and constrain a KL update.

\subsection{Trust-region policy optimization}
\label{subsec:rl-trust-region-policy-optimization}

The natural gradient supplies a geometrically meaningful direction, but a
direction alone does not specify a safe step length. A large update can leave
the region in which a local approximation is reliable. Trust Region Policy
Optimization (TRPO) addresses this issue by maximizing a surrogate objective
while constraining the divergence between the old and new policies
\parencite{schulman2015trpo}.

For a current policy $\pi$, the performance difference identity in
\cref{eq:rl-performance-difference} depends on the state-visitation frequency
$\rho_{\widetilde\pi}$ of the new policy. Directly optimizing this expression
is difficult because the state distribution changes with the policy. TRPO
replaces this new-policy occupancy with the old-policy occupancy and defines
\begin{equation}
    L_\pi(\widetilde\pi)
    =\eta(\pi)
     +\sum_s\rho_\pi(s)
       \sum_a\widetilde\pi(a\mid s)A^\pi(s,a).
    \label{eq:rl-trpo-surrogate}
\end{equation}
At the old policy, $L_\pi(\pi)=\eta(\pi)$ because the expected advantage is
zero. For a differentiable parameterized policy, the surrogate and the true
objective also have the same first derivative at the old parameter value.
This makes the surrogate a first-order local approximation, but it does not
justify arbitrarily large policy changes.

To bound the error caused by the changed state distribution, define total
variation divergence for discrete distributions by
\begin{equation}
    D_{\mathrm{TV}}(p,q)
    =\frac{1}{2}\sum_a|p(a)-q(a)|,
    \qquad
    \alpha=\max_sD_{\mathrm{TV}}
    (\pi(\cdot\mid s),\widetilde\pi(\cdot\mid s)),
    \label{eq:rl-total-variation}
\end{equation}
and let $\epsilon=\max_{s,a}|A^\pi(s,a)|$. The policy-improvement theorem
proved by \textcite{schulman2015trpo} states
\begin{equation}
    \eta(\widetilde\pi)
    \geq L_\pi(\widetilde\pi)
       -\frac{4\epsilon\gamma}{(1-\gamma)^2}\alpha^2.
    \label{eq:rl-trpo-tv-bound}
\end{equation}
The quadratic dependence on $\alpha$ reflects the fact that two policies can
be coupled to select the same action with high probability when their
state-wise total variation distance is small. Their trajectories remain
identical until a disagreement occurs; the resulting occupancy mismatch is
therefore controlled by a second-order term.

Using the inequality
$D_{\mathrm{TV}}(p,q)^2\leq D_{\mathrm{KL}}(p\|q)$, define
\begin{equation}
    D_{\mathrm{KL}}^{\max}(\pi,\widetilde\pi)
    =\max_sD_{\mathrm{KL}}
      (\pi(\cdot\mid s)\|\widetilde\pi(\cdot\mid s)),
    \qquad
    C=\frac{4\epsilon\gamma}{(1-\gamma)^2}.
    \label{eq:rl-trpo-kl-definitions}
\end{equation}
Then the lower bound becomes
\begin{equation}
    \eta(\widetilde\pi)
    \geq L_\pi(\widetilde\pi)
       -C D_{\mathrm{KL}}^{\max}(\pi,\widetilde\pi).
    \label{eq:rl-trpo-kl-bound}
\end{equation}
The maximum state-wise divergence is the quantity required by this theoretical
bound. Maximizing the right-hand side at every iteration produces a
non-decreasing sequence of true policy performances under the theorem's
assumptions. In particular, the current policy is feasible and gives equality
in the lower bound, so maximizing the lower-bound objective cannot reduce its
value.

The exact maximum-KL constraint is difficult to enforce over a large state
space. TRPO therefore uses a practical average-divergence approximation. For a
state distribution $\rho$, define
\begin{equation}
    D_{\mathrm{KL}}^{\rho}(\theta_{\mathrm{old}},\theta)
    =\mathbb{E}_{s\sim\rho}
      \left[
      D_{\mathrm{KL}}\bigl(
      \pi_{\theta_{\mathrm{old}}}(\cdot\mid s)
      \|\pi_\theta(\cdot\mid s)\bigr)
      \right].
    \label{eq:rl-trpo-average-kl}
\end{equation}
The practical constrained problem is
\begin{equation}
    \begin{aligned}
        \max_\theta\quad&L_{\theta_{\mathrm{old}}}(\theta)\\
        \text{subject to}\quad&
        D_{\mathrm{KL}}^{\rho_{\mathrm{old}}}
        (\theta_{\mathrm{old}},\theta)\leq\delta.
    \end{aligned}
    \label{eq:rl-trpo-practical-problem}
\end{equation}
Replacing the maximum state-wise constraint by an average is a practical
approximation; it does not, without additional assumptions, retain the exact
monotonic-improvement guarantee of \cref{eq:rl-trpo-kl-bound}. This distinction
between the theorem and the implementable approximation is essential when
interpreting TRPO.

The objective can be estimated from trajectories sampled under the old policy.
An importance-sampling representation of the action expectation is
\begin{equation}
    \mathbb{E}_{a\sim\pi_{\mathrm{old}}(\cdot\mid s)}
    \left[
    \frac{\pi_\theta(a\mid s)}
         {\pi_{\mathrm{old}}(a\mid s)}
    A^{\pi_{\mathrm{old}}}(s,a)
    \right],
    \label{eq:rl-trpo-importance-ratio}
\end{equation}
provided the sampling distribution has support wherever the updated policy
assigns probability. TRPO describes a single-path estimator, which samples
trajectories from the old policy and estimates action values from discounted
rewards along each trajectory. It also describes a vine estimator, which
selects states from sampled trajectories, evaluates multiple actions from each
selected state, and can use common random numbers across rollouts to reduce the
variance of action-value comparisons. The single-path method is applicable in
model-free settings where the system cannot be reset to an arbitrary state; the
vine method is particularly natural in a simulator that permits such resets
\parencite{schulman2015trpo}.

The local constrained problem connects TRPO directly to the natural gradient.
Let $g=\nabla_\theta L_{\theta_{\mathrm{old}}}(\theta)$ evaluated at
$\theta_{\mathrm{old}}$. A second-order expansion of the KL constraint gives
\begin{equation}
    D_{\mathrm{KL}}^{\rho_{\mathrm{old}}}
    (\theta_{\mathrm{old}},\theta_{\mathrm{old}}+\Delta\theta)
    \simeq\frac{1}{2}\Delta\theta^{\mathsf T}F\Delta\theta,
    \label{eq:rl-trpo-fisher-quadratic}
\end{equation}
where $F$ is the Fisher information matrix under the old-policy state
weighting. The local optimization is therefore
\begin{equation}
    \max_{\Delta\theta}g^{\mathsf T}\Delta\theta
    \quad\text{subject to}\quad
    \frac{1}{2}\Delta\theta^{\mathsf T}F\Delta\theta\leq\delta.
    \label{eq:rl-trpo-local-problem}
\end{equation}
The Lagrange-multiplier solution is proportional to $F^{-1}g$. If
$s=F^{-1}g$, the step length that saturates the quadratic constraint is
\begin{equation}
    \beta=\sqrt{\frac{2\delta}{s^{\mathsf T}Fs}},
    \qquad
    \Delta\theta=\beta s.
    \label{eq:rl-trpo-step-length}
\end{equation}
Thus, the Fisher matrix determines the direction and the trust-region radius
determines the scale.

For a neural policy with $d$ parameters, explicitly storing a dense Fisher
matrix requires $O(d^2)$ memory, and a dense factorization or inverse can
require $O(d^3)$ arithmetic operations. TRPO avoids forming the dense inverse
by solving $Fs=g$ approximately with conjugate gradients. Conjugate gradients
requires repeated matrix--vector products $Fv$, which can be obtained through
Hessian--vector products of the KL objective using automatic differentiation.
The number of conjugate-gradient iterations controls the accuracy of the
linear-system solution. After obtaining a candidate direction, a backtracking
line search evaluates the actual nonlinear surrogate and KL divergence; it
reduces the step until the candidate improves the surrogate and satisfies the
practical divergence constraint. The line search is needed because the
first-order objective approximation and quadratic KL approximation are local,
not globally exact.

\begin{algorithm}[H]
\caption{Sample-based trust-region policy optimization}
\label{alg:rl-trpo}
\KwIn{Initial policy parameters $\theta_0$, KL radius $\delta$,
      trajectory sampler, conjugate-gradient tolerance, and line-search rule.}
\KwOut{Updated policy parameters $\theta$.}
Set $\theta\leftarrow\theta_0$\;
\While{the stopping condition is not satisfied}{
    Set $\theta_{\mathrm{old}}\leftarrow\theta$\;
    Collect trajectories using $\pi_{\theta_{\mathrm{old}}}$ and estimate advantages\;
    Estimate the surrogate gradient $g$ and Fisher-vector products $v\mapsto Fv$\;
    Approximately solve $Fs=g$ with conjugate gradients\;
    Scale $s$ using the quadratic KL constraint to obtain a candidate step\;
    Evaluate backtracking candidates using the sampled surrogate and average KL\;
    \If{a candidate improves the surrogate and satisfies the KL constraint}{
        Accept the candidate as the new $\theta$\;
    }
    \Else{
        Retain $\theta_{\mathrm{old}}$ and terminate if the line search fails\;
    }
}
\end{algorithm}

The algorithm in \cref{alg:rl-trpo} separates statistical estimation from
numerical optimization. Its sample quality depends on the trajectory estimator
and advantage estimates; its local optimization quality depends on the
Fisher-vector products and conjugate-gradient accuracy; and its final accepted
step depends on the actual surrogate and KL evaluations. TRPO is therefore not
an exact solution of the original constrained problem, but a practical
approximation to a theoretically motivated policy-improvement procedure
\parencite{schulman2015trpo}.

\begin{figure}[htbp]
    \centering
    \resizebox{\textwidth}{!}{% Native LaTeX picture: sample estimation and trust-region numerical update.
\setlength{\unitlength}{1mm}
\begin{picture}(160,55)
    \put(2,34){\framebox(38,14){\shortstack{Old policy\\$\pi_{\theta_{\rm old}}$}}}
    \put(61,34){\framebox(38,14){\shortstack{Sample trajectories\\and estimate $A$}}}
    \put(120,34){\framebox(38,14){\shortstack{Surrogate gradient\\$g=\nabla_\theta L$}}}
    \put(40,41){\vector(1,0){21}}
    \put(99,41){\vector(1,0){21}}
    \put(120,6){\framebox(38,14){\shortstack{Line search\\accept / reject}}}
    \put(61,6){\framebox(38,14){\shortstack{Conjugate gradients\\solve $Fs=g$}}}
    \put(2,6){\framebox(38,14){\shortstack{KL constraint\\$D_{\rm KL}\leq\delta$}}}
    \put(139,34){\vector(0,-1){14}}
    \put(120,13){\vector(-1,0){21}}
    \put(61,13){\vector(-1,0){21}}
\end{picture}
}
    \caption{The sample-based TRPO pipeline, from old-policy trajectories and advantage
    estimates through the Fisher-system solve to a line-searched update.}
    \label{fig:rl-trpo-geometry}
\end{figure}
% Alt text for accessibility register: The policy parameterization induces a
% trajectory distribution; its Fisher matrix scales the gradient, and a KL
% radius limits the resulting parameter update.

\subsection{Regularized Markov decision processes}
\label{subsec:rl-regularized-mdps}

Trust-region optimization limits the change between successive policies. A
related but more general strategy modifies the policy optimization objective
through a regularizer. The theory developed by
\textcite{geist2019regularized} treats regularized Markov decision processes
using convex analysis and a generalized Bellman framework. The central idea is
to define the regularizer over action distributions at each state and then
incorporate it into the Bellman evaluation and optimality operators.

Let $\Omega:\Delta_{\mathcal A}\rightarrow\mathbb{R}$ be a strongly convex
regularizer. For a value function $v$, define the one-step action-value vector
\begin{equation}
    q_v(s,a)
    =r(s,a)+\gamma\sum_{s'}P(s'\mid s,a)v(s').
    \label{eq:rl-one-step-q}
\end{equation}
The regularized greedy distribution at state $s$ solves
\begin{equation}
    \pi^*_s
    =\operatorname*{arg\,max}_{p\in\Delta_{\mathcal A}}
      \left[\langle p,q_v(s,\cdot)\rangle-\Omega(p)\right].
    \label{eq:rl-regularized-greedy}
\end{equation}
The linear term favors actions with high expected return; the regularizer
changes the preferred distribution according to the chosen convex geometry.
Strong convexity makes the maximizer unique under the framework's assumptions.

The Legendre--Fenchel convex conjugate of $\Omega$ is
\begin{equation}
    \Omega^*(q)
    =\max_{p\in\Delta_{\mathcal A}}
      \left[\langle p,q\rangle-\Omega(p)\right].
    \label{eq:rl-legendre-fenchel}
\end{equation}
When the conjugate is differentiable, its gradient returns the unique
maximizing distribution:
\begin{equation}
    \pi_s^*=\nabla\Omega^*(q_v(s,\cdot)).
    \label{eq:rl-conjugate-policy}
\end{equation}
The convex conjugate therefore converts a regularized optimization over the
probability simplex into a value function of the action-value vector.

For negative Shannon entropy,
\begin{equation}
    \Omega(p)=\sum_a p(a)\log p(a),
    \label{eq:rl-negative-entropy}
\end{equation}
with the convention $0\log0=0$, the conjugate is the log-sum-exp function:
\begin{equation}
    \Omega^*(q)=\log\sum_a e^{q(a)},
    \qquad
    \nabla\Omega^*(q)(a)
    =\frac{e^{q(a)}}{\sum_b e^{q(b)}}.
    \label{eq:rl-softmax-conjugate}
\end{equation}
Thus, negative-entropy regularization yields a softmax policy. If the regularizer
is scaled by a temperature $\tau>0$, so that $\Omega_\tau(p)=
\tau\sum_a p(a)\log p(a)$, its conjugate is
$\Omega_\tau^*(q)=\tau\log\sum_a e^{q(a)/\tau}$ and the policy is
$\pi(a)=e^{q(a)/\tau}/\sum_b e^{q(b)/\tau}$. For example, the quadratic regularizer
$\Omega(p)=\tfrac{1}{2}(\|p\|_2^2-1)$ is associated with the order-two
Tsallis construction and yields a sparsemax-type solution, which can assign
exactly zero probability to some actions. These examples show that the regularizer
 determines the form of the regularized greedy policy
\parencite{geist2019regularized}.

The regularized Bellman evaluation operator is
\begin{equation}
    [T_{\pi,\Omega}v](s)
    =\langle\pi(\cdot\mid s),q_v(s,\cdot)\rangle
     -\Omega(\pi(\cdot\mid s)),
    \label{eq:rl-regularized-evaluation}
\end{equation}
and the regularized optimality operator is
\begin{equation}
    [T^*_{\Omega}v](s)
    =\max_{p\in\Delta_{\mathcal A}}
      \left[\langle p,q_v(s,\cdot)\rangle-\Omega(p)\right]
    =\Omega^*(q_v(s,\cdot)).
    \label{eq:rl-regularized-optimality}
\end{equation}
The regularized value functions are fixed points of these operators:
$V^{\pi,\Omega}=T_{\pi,\Omega}V^{\pi,\Omega}$ and
$V^{*,\Omega}=T^*_{\Omega}V^{*,\Omega}$. The regularized greedy policy is
$G_\Omega(v)=\nabla\Omega^*(q_v)$.

The regularized operators preserve the key properties needed for dynamic
programming. They are monotone: if $v_1\geq v_2$ componentwise, then
$T_{\pi,\Omega}v_1\geq T_{\pi,\Omega}v_2$ and
$T^*_{\Omega}v_1\geq T^*_{\Omega}v_2$. They also satisfy constant-shift
identities,
\begin{equation}
    T_{\pi,\Omega}(v+c\mathbf{1})
    =T_{\pi,\Omega}v+\gamma c\mathbf{1},
    \qquad
    T^*_{\Omega}(v+c\mathbf{1})
    =T^*_{\Omega}v+\gamma c\mathbf{1},
    \label{eq:rl-regularized-shift}
\end{equation}
because transition probabilities sum to one and the regularizer does not
depend on $v$. Finally, both operators are $\gamma$-contractions in the
supremum norm:
\begin{equation}
    \|T_{\pi,\Omega}v_1-T_{\pi,\Omega}v_2\|_\infty
    \leq\gamma\|v_1-v_2\|_\infty,
    \qquad
    \|T^*_{\Omega}v_1-T^*_{\Omega}v_2\|_\infty
    \leq\gamma\|v_1-v_2\|_\infty.
    \label{eq:rl-regularized-contraction}
\end{equation}
The regularized Bellman operators therefore have unique fixed points, and
regularized value iteration converges under the same discounted contraction
logic as its unregularized counterpart \parencite{geist2019regularized}.

A fixed regularizer changes the objective, so the optimal regularized policy
need not be optimal for the original unregularized MDP. If the regularizer is
bounded on the simplex,
\begin{equation}
    L_\Omega\leq\Omega(p)\leq U_\Omega
    \qquad\text{for all }p\in\Delta_{\mathcal A},
    \label{eq:rl-regularizer-bounds}
\end{equation}
then the regularized and unregularized values differ by a bounded amount.
In particular, the unregularized value of the regularized optimal policy
satisfies
\begin{equation}
    V^* - V^{\pi^{*,\Omega}}
    \leq\frac{U_\Omega-L_\Omega}{1-\gamma}\mathbf{1}.
    \label{eq:rl-regularization-gap}
\end{equation}
The range of the regularizer and the discount horizon therefore determine a
bound on the cost of solving the regularized rather than the original problem.
This is distinct from a claim that regularization always improves the
unregularized optimum; it controls the difference between the two objectives
under the stated boundedness assumptions \parencite{geist2019regularized}.

\begin{figure}[htbp]
    \centering
    \resizebox{\textwidth}{!}{% Native LaTeX picture: regularized greedy map and proximal update.
\setlength{\unitlength}{1mm}
\begin{picture}(160,58)
    \put(2,38){\framebox(40,13){\shortstack{Action values\\$q_v(s,\cdot)$}}}
    \put(60,38){\framebox(40,13){\shortstack{Convex conjugate\\$\Omega^*(q)$}}}
    \put(118,38){\framebox(40,13){\shortstack{Regularized policy\\$\nabla\Omega^*(q)$}}}
    \put(42,44){\vector(1,0){18}}
    \put(100,44){\vector(1,0){18}}
    \put(2,8){\framebox(40,13){\shortstack{Previous policy\\$\pi_k$}}}
    \put(60,8){\framebox(40,13){\shortstack{Bregman divergence\\$D_\Omega(\pi\|\pi_k)$}}}
    \put(118,8){\framebox(40,13){\shortstack{Proximal update\\$\arg\max_\pi\{\langle q,\pi\rangle-D\}$}}}
    \put(42,14){\vector(1,0){18}}
    \put(100,14){\vector(1,0){18}}
\end{picture}
}
    \caption{Regularized policy improvement. The action-value vector is passed
    through the convex conjugate to obtain a regularized policy; in the
    proximal form, the new policy is additionally kept close to the previous
    policy by a Bregman divergence.}
    \label{fig:rl-regularized-mdp}
\end{figure}
% Alt text for accessibility register: Action values feed a convex-conjugate
% mapping to a regularized policy. A second path shows the previous policy and
% Bregman divergence entering a proximal policy update.

\subsection{Mirror descent and modified policy iteration}
\label{subsec:rl-mirror-descent-mdpi}

A fixed regularizer modifies the optimal solution of the MDP. An alternative
is to regularize each update by penalizing the divergence between the new policy
and the policy from the previous iteration. This makes the regularization
iteration-dependent and connects policy improvement to proximal convex
optimization.

For a differentiable convex function $\Omega$, the Bregman divergence from
$q$ to $p$ is
\begin{equation}
    D_\Omega(p\|q)
    =\Omega(p)-\Omega(q)
     -\langle\nabla\Omega(q),p-q\rangle.
    \label{eq:rl-bregman-divergence}
\end{equation}
Convexity implies $D_\Omega(p\|q)\geq0$ and
$D_\Omega(q\|q)=0$. In general, it is not symmetric and therefore is not a
metric. For negative Shannon entropy, the Bregman divergence is the KL
divergence:
\begin{equation}
    D_\Omega(p\|q)
    =\sum_a p(a)\log\frac{p(a)}{q(a)}
    =D_{\mathrm{KL}}(p\|q).
    \label{eq:rl-bregman-kl}
\end{equation}

The proximal, mirror-descent policy update is
\begin{equation}
    \pi_{k+1}
    =\operatorname*{arg\,max}_{\pi\in\Delta_{\mathcal A}}
      \left[\langle q_k,\pi\rangle
      -D_\Omega(\pi\|\pi_k)\right].
    \label{eq:rl-mirror-policy-update}
\end{equation}
The negative action-value vector plays the role of a gradient in the usual
minimization form of mirror descent. The Bregman term controls the change from
the previous policy in the geometry induced by $\Omega$.

For the KL case, the update can be derived explicitly. Substitute
$D_{\mathrm{KL}}(\pi\|\pi_k)$ into \cref{eq:rl-mirror-policy-update}.
Introducing a Lagrange multiplier for the normalization constraint
$\sum_a\pi(a)=1$ gives the stationarity condition
\begin{equation}
    q_k(a)-\log\frac{\pi_{k+1}(a)}{\pi_k(a)}-1+\lambda=0.
    \label{eq:rl-kl-mirror-stationarity}
\end{equation}
Solving and normalizing yields
\begin{equation}
    \pi_{k+1}(a)
    =\frac{\pi_k(a)e^{q_k(a)}}
           {\sum_b\pi_k(b)e^{q_k(b)}}.
    \label{eq:rl-kl-mirror-update}
\end{equation}
The old policy is therefore reweighted by exponentiated action values. Unlike
an independently recomputed softmax, this update retains the previous policy
as its reference distribution.

Mirror Descent Modified Policy Iteration (MD-MPI) combines this proximal
greedy step with partial policy evaluation. Its two forms differ in the
evaluation operator:
\begin{equation}
    \begin{aligned}
        \text{Type 1:}\quad
        &\pi_{k+1}=G_{\Omega^{\pi_k}}(v_k),\\
        &v_{k+1}=(T_{\pi_{k+1},\Omega^{\pi_k}})^m v_k;
    \end{aligned}
    \qquad
    \begin{aligned}
        \text{Type 2:}\quad
        &\pi_{k+1}=G_{\Omega^{\pi_k}}(v_k),\\
        &v_{k+1}=(T_{\pi_{k+1}})^m v_k.
    \end{aligned}
    \label{eq:rl-mdmpi-types}
\end{equation}
Both types regularize the greedy step relative to the previous policy. Type 1
also uses the corresponding regularized evaluation operator; Type 2 uses the
ordinary evaluation operator. The distinction changes where approximation
errors enter the iteration and consequently changes the details of the error
analysis.

The framework provides a useful unification. In the formulation discussed by
\textcite{geist2019regularized}, Type 2 with complete policy evaluation and a
KL divergence recovers the structure of TRPO, up to the practical distinction
between a hard divergence constraint and a divergence penalty in the proximal
objective. The same framework relates other algorithms to regularized
modified policy iteration, but those algorithms are not developed further
here because the present section is restricted to the four supplied papers.

The analysis of MD-MPI introduces a cumulative regret quantity
\begin{equation}
    L_K=\sum_{k=1}^{K}(V^*-V^{\pi_k})
    \label{eq:rl-mdmpi-regret}
\end{equation}
and the initial Bregman radius
\begin{equation}
    R_{\Omega}^{\pi_0}
    =\sup_{\pi}D_\Omega(\pi\|\pi_0),
    \label{eq:rl-mdmpi-radius}
\end{equation}
with the appropriate state-wise norm convention. The key telescoping term is
\begin{equation}
    \begin{aligned}
    \sum_{k=0}^{K-1}
    \bigl[D_\Omega(\pi^*\|\pi_k)
         -D_\Omega(\pi^*\|\pi_{k+1})\bigr]
    &=D_\Omega(\pi^*\|\pi_0)
      -D_\Omega(\pi^*\|\pi_K)\\
    &\leq D_\Omega(\pi^*\|\pi_0).
    \end{aligned}
    \label{eq:rl-bregman-telescoping}
\end{equation}
This is the mechanism by which the initial divergence controls the cumulative
contribution of the proximal updates. With a uniform initial policy and KL
divergence over a finite action set,
\begin{equation}
    \sup_{\pi\in\Delta_{\mathcal A}}
    D_{\mathrm{KL}}(\pi\|\pi_0)=\log|\mathcal A|.
    \label{eq:rl-uniform-kl-radius}
\end{equation}
The logarithmic dependence on the number of actions follows because the
maximum is attained by a distribution concentrated on one action.

In the exact setting, the regret bound in the MD-MPI analysis implies a bound
of the form
\begin{equation}
    \frac{1}{K}\|L_K\|_\infty
    \leq
    \frac{1-\gamma^K}{K(1-\gamma)^2}
       2\gamma\|V^*-v_0\|_\infty
    +\frac{R_{\Omega}^{\pi_0}}{K}.
    \label{eq:rl-mdmpi-exact-regret-bound}
\end{equation}
Approximate versions add terms for the greedy and evaluation errors propagated
through policy transition kernels. The regret bound controls the average loss
across iterates; it does not necessarily say that the last iterate is the best
one. The analysis consequently gives a bound on the best policy generated so
far through
\begin{equation}
    \min_{1\leq k\leq K}
    \|V^*-V^{\pi_k}\|_{1,\rho}
    \leq\frac{1}{K}\|L_K\|_{p,\rho},
    \label{eq:rl-best-iterate-regret}
\end{equation}
for the norm and distribution conditions specified in the paper. This is why
the best iterate can be a more appropriate output than the final iterate
\parencite{geist2019regularized}.

\subsection{Worked example and connection to autoregressive language models}
\label{subsec:rl-worked-example-language-models}

A two-action calculation makes the KL-proximal update concrete. Consider one
state with actions $a_1$ and $a_2$, previous policy
$\pi_k=(0.5,0.5)$, and action-value vector $q_k=(1,0)$. Under the KL mirror
update in \cref{eq:rl-kl-mirror-update},
\begin{equation}
    \pi_{k+1}(a_1)
    =\frac{0.5e^1}{0.5e^1+0.5e^0}
    =\frac{e}{e+1}\simeq0.7311,
    \qquad
    \pi_{k+1}(a_2)\simeq0.2689.
    \label{eq:rl-numerical-mirror-update}
\end{equation}
The policy assigns more probability to the action with the larger value, but it
retains nonzero probability for the other action. The KL divergence from the
updated policy to the previous uniform policy is
\begin{equation}
    \begin{aligned}
    D_{\mathrm{KL}}(\pi_{k+1}\|\pi_k)
    &\simeq0.7311\log\frac{0.7311}{0.5}
      +0.2689\log\frac{0.2689}{0.5}\\
    &\simeq0.1109\ \text{nats}.
    \end{aligned}
    \label{eq:rl-numerical-kl}
\end{equation}
The unit is nats because the logarithm is natural. This example illustrates the
mechanism of the proximal update: action values shift the distribution, while
the divergence is measured relative to the previous policy. A larger action
value difference produces a more concentrated update; a stronger divergence
penalty, if included as a coefficient, would make the update more conservative.

For an autoregressive language model, a natural MDP representation takes the
generated token prefix as the state and the next token as the action. If the
current prefix is $s_t=(x_1,\ldots,x_t)$, the action is
$a_t=x_{t+1}$, and the transition appends the selected token:
\begin{equation}
    s_{t+1}=s_t\mathbin{\Vert}a_t,
    \qquad
    \pi_\theta(a_t\mid s_t)
    =p_\theta(x_{t+1}\mid x_1,\ldots,x_t).
    \label{eq:rl-language-model-mdp}
\end{equation}
Here $\mathbin{\Vert}$ denotes sequence concatenation. If the state contains
the complete prefix and any other context needed to determine the reward and
transition, this representation is Markov: the next state is determined by
the current prefix and selected token. A reward may be assigned at each token
step or at the end of the generated sequence. With a terminal sequence reward,
the intermediate rewards can be zero and the final transition carries the
sequence-level outcome. The horizon is the maximum generated length or a
stopping time determined by the generation process.

This mapping is an application of the MDP definitions developed above; the
four cited theoretical papers do not themselves establish language-model-
specific reinforcement-learning results. Its value for the present book is
conceptual and mathematical. It makes explicit why token generation can be
viewed as sequential decision-making, why a sequence-level reward creates a
delayed-credit-assignment problem, and why policy optimization changes a
distribution over complete generated sequences rather than merely changing an
isolated token prediction. The policy is the language model's conditional next-
token distribution, the trajectory is a generated sequence, and the return is
the discounted or finite-horizon accumulation of rewards associated with that
sequence.

\begin{figure}[htbp]
    \centering
    \resizebox{\textwidth}{!}{% Native LaTeX picture: autoregressive token generation as an MDP trajectory.
\setlength{\unitlength}{1mm}
\begin{picture}(170,48)
    \put(1,20){\framebox(39,17){\shortstack{State\\prefix $s_t=(x_1,\ldots,x_t)$}}}
    \put(47,20){\framebox(34,17){\shortstack{Action\\next token $a_t$}}}
    \put(88,20){\framebox(39,17){\shortstack{Transition\\append $a_t$}}}
    \put(134,20){\framebox(35,17){\shortstack{Reward\\token or sequence}}}
    \put(40,28){\vector(1,0){7}}
    \put(81,28){\vector(1,0){7}}
    \put(127,28){\vector(1,0){7}}
    \put(43,11){\makebox(85,5){$s_{t+1}=s_t\mathbin{\Vert}a_t$}}
\end{picture}
}
    \caption{A natural MDP interpretation of autoregressive generation. The
    current token prefix is the state, the next token is the action, appending
    the token produces the next state, and a reward may be attached to a token
    step or to the completed sequence.}
    \label{fig:rl-language-model-mdp}
\end{figure}
% Alt text for accessibility register: Four boxes show a token prefix as state,
% a next-token choice as action, a deterministic append operation as transition,
% and a token-level or sequence-level reward.

\subsection{Limitations and connection to the next topic}
\label{subsec:rl-theory-limitations-next-topic}

The results in this section rely on assumptions that must remain visible when
they are applied to large neural policies. The contraction arguments use a
discount factor strictly below one and bounded rewards. The finite-dimensional
matrix derivations assume finite state and action spaces, although the
operator formulation extends to probability kernels. Exact policy-improvement
claims assume accurate value functions and greedy updates. Approximate methods
replace those operations with estimates, so their guarantees depend on both
the approximation errors and the state-distribution coefficients that
propagate those errors.

The natural-gradient construction depends on the chosen trajectory
probability model and on the regularity of the policy parameterization. A
singular Fisher matrix indicates redundant or locally unidentifiable parameter
directions. TRPO's monotonic-improvement theorem uses a maximum state-wise KL
constraint, whereas the practical algorithm uses an average KL constraint and
sample-based estimates. Its practical stability should therefore not be
confused with an unconditional theorem for every neural policy and sampling
regime. Fixed regularization also changes the optimization objective, while
Bregman regularization controls updates relative to a previous policy and
introduces its own dependence on the initial divergence and approximation
errors.

Taken together, the four supplied papers establish a progression from exact
dynamic programming to approximate policy iteration, from coordinate-dependent
gradients to information-geometric updates, and from unconstrained policy
improvement to trust-region and Bregman-regularized optimization. The common
mathematical structure is the control of policy change under imperfect
estimation. For language models, this structure explains how next-token
probability distributions can be optimized against sequence-level feedback
without treating each generated token as an independent supervised target.
The next section can now introduce actor--critic architectures, in which a
policy model and a value estimator are trained together to estimate the
advantages needed for policy improvement.
